
\subsection{Zitieren}
\label{sec:zitieren}

\begin{itemize}
\item Wenn Sie ein Dokument verfassen, das in irgendeiner Weise verbreitet wird (oder werden könnte), sollten Sie darauf achten, dass alle Ihre Abbildungen frei von Urheberrechtsproblemen sind. Das bedeutet: Entweder Sie zeichnen sie selbst (auch die Übernahme von Teilen aus anderen Bildern und deren Veränderung ist ein Problem) oder Sie holen eine Genehmigung ein und geben in der Bildunterschrift an, dass Sie diese Genehmigung erhalten haben und von wem. Meistens sagen Ihnen die Verlage, wie man das schreibt, und haben sogar ein automatisiertes Verfahren, um die Erlaubnis zu erhalten, z. B. über RightsLink/Copyright Clearance Center (meist sofort und kostenlos für Abschlussarbeiten von Studierenden). Auch Bilder aus dem öffentlichen Bereich müssen ordnungsgemäß referenziert werden. Es gibt auch verschiedene Arten von Lizenzen, z. B. erlauben einige keine Änderungen an einem Bild.

\item Vermeiden Sie direkte Zitate (``...'') wann immer möglich, formulieren Sie stattdessen Aussagen um, wenn Sie zitieren.

\item Es ist eindeutig vorzuziehen, Primär- und Originalliteratur zu zitieren, nicht Rezensionen oder Lehrbücher. Gehen Sie auf die Originalquellen zurück. Wenn Sie Rezensionen oder Bücher zitieren, tun Sie dies ausdrücklich, z. B. mit ``ein Überblick findet sich in...''.

\item  \item Lesen Sie jede Arbeit, die Sie zitieren, so lange, bis Sie sie verstehen, oder zumindest den Teil der Arbeit, auf den Sie sich beziehen. Verlassen Sie sich nicht nur auf Zusammenfassungen oder, schlimmer noch, auf Sekundärliteratur, die die Originalquelle (oft falsch) zitiert. Wenn Sie es für unmöglich halten, die zitierten Quellen zu lesen und zu verstehen, schränken Sie den Umfang Ihrer Arbeit so weit ein, bis es möglich ist.

\item Belegen Sie jede einzelne Behauptung, die keine offensichtliche Tatsache ist, entweder durch einen Verweis oder durch ein eigenes Experiment, eine Schlussfolgerung oder eine Berechnung. 

\cite{Menge}
\end{itemize}